% Folder 156-9, batch 03 (archivists' pages 41-60; author's pp. 38-57).
% Pass: Opus 5.5 (claude-opus-5-5), 2026-10-03 — first pass, unchecked against the pages by a human.
%
% Notation fixed for this batch: \mathcal{A} for the atelier (his curly A),
% \trianglelefteq for his "sous-figure" relation, \mathrel{|\circ|} for the
% disjunction relation, \overset{\circ}{\ll} for his circled «, \complement for
% his C (complement in a spatios), Dép for « déploiement », X° for the support
% of X, \widetilde{X} for the set of sub-figures of X. "Atspat", "Atens",
% "Figél", "Figspat" are his abbreviations (atelier spatial, atelier
% ensembliste, figures élémentaires, figures spatiales).
\documentclass[11pt,a4paper]{article}
\input{../preamble/grothendieck}

\folder{156-9}
\batch{3}
\pages{41}{60}
\foldertitle{[Chapitre] IX et IX bis. [Ateliers] : notes manuscrites (05-15/07/1986).}
\dating{1986}
\watermark{Édition de démonstration}

\begin{document}

\page{41}
\note{p. 38 de l'auteur. La phrase commence sur la page précédente, hors du lot.}

… y jouent \uncertain{dans le rôle} des Inf quelconques, contenant
$\varnothing$ (plus petit élément) et $\mathcal{A}$ (plus grand élément). On y
a une opération $\operatorname{cosupp}^{\circ}$ \uncertain{ou} $\complement$,
qui est une anti-involution. Sa \ill{} \struck{implique} (\uncertain{i.e.}
$\complement$) implique l'existence des Sup quelconques dans
$\Sigma_{\mathcal{M}}$ \uncertain{(anti-Sup)}, et $\complement$ échange Inf et
Sup, plus petit et plus grand élément. La relation $\mathrel{|\circ|}$
\struck{se transporte pour} \add{s'étend} aux éléments de
$\Sigma_{\mathcal{M}}$, par
\note{l'indice de $\Sigma$, ici et deux lignes plus haut, est une lettre qui
ressemble à un $\mathcal{M}$ ; on attendrait $\Sigma_{\mathcal{A}}$, qui
figure en toutes lettres au bas de la page.}
\[
  S \mathrel{|\circ|} S' \iff \forall X \in S,\ X' \in S',\ \text{on a }
  X \mathrel{|\circ|} X' .
\]
\note{au-dessus de la double flèche, un mot inséré, \uncertain{par définition}.}

On a
\[
  S \mathrel{|\circ|} S' \iff S \subset \complement S' \iff S' \subset \complement S .
\]
On a
\[
  S \mathrel{|\circ|} S' \Longrightarrow S \cap S' = \varnothing ,
\]
l'inverse n'est pas vrai en général.

\marginal{(en diagonale, à travers les formules) \uncertain{On a aussi} \ill{}
\uncertain{plus petit et plus grand} \ill{}}

La donnée d'un \ill{} \ill{}, soit $\mathfrak{F}$, (\uncertain{comme}
« figures »), \struck{permet} plus la donnée de $\trianglelefteq$, permet de
définir dans $\Sigma_{\mathcal{A}}$ une partie, formée des \emph{supports
constructibles}

\marginal{Structure d'un spatios : a) relation d'ordre b) application
$\complement$. \uncertain{Axiomes} : 1°) $\exists$ plus petit et plus grand
élément ; 2°) $\complement$ est une anti-involution ; 3°) \struck{$\forall x$}
$x \wedge \complement x = \varnothing_{\Sigma}$ ($\Leftrightarrow$
$x \vee \complement x = 1_{\Sigma}$). Description de $\mathcal{A}$
\uncertain{dans} $\operatorname{Fig\acute{e}l\,ord}(\Sigma)$.}

\page{42}
\note{p. 39 de l'auteur.}

\textbf{$\Sigma\mathrm{cons}_{\mathcal{A}}$}

Elle n'est \struck{\ill{}} en général stable ni par Sup, ni par Inf, ni par
$\complement$. Mais il y a un axiome (« axiome des supports ») qui tient lieu
d'une telle compatibilité, sur lequel nous allons \add{encore} revenir --- et à
un \uncertain{moment} \add{\uncertain{cet}} axiome, en faisant intervenir cette
fois aussi la structure $\trianglelefteq$.

\struck{Pour l'instant, je suis, \ill{}} \add{Du point de vue des notations,}
pour $X \in \mathcal{A}$, je dénote par
\[
  X^{\circ} \in \Sigma\mathrm{cons}_{\mathcal{A}}
\]
l'élément \struck{\ill{}} \add{$\operatorname{supp}^{\circ}(X)$}
$\operatorname{supp}^{\circ} X$, qui peut être décrit comme le plus petit
$S \in \Sigma_{\mathcal{A}}$ (élément aussi de
$\Sigma\mathrm{cons}_{\mathcal{A}}$) tel que $X \in S$. C'est une
explicitation possible de l'intuition de la \emph{\uncertain{partie occupée}}
\uncertain{associée} à une \uncertain{figure} donnée $X$.

\page{43}
\note{p. 40 de l'auteur.}

Je vais maintenant relier la relation \add{(\uncertain{déjà}
$\overset{\circ}{\ll}$)} $\overset{\circ}{\ll}$, qui n'est pas encore
intervenue dans les considérations de supports, \ill{} ce qui \uncertain{ne
veut} pas dire que les $S \in \Sigma_{\mathcal{A}}$ sont des parties
\uncertain{$\mathfrak{F}$-fermées} de $\mathcal{A}$), à ces notions de
support.

Tout d'abord, pour tout $X \in \mathcal{A}$, considérons \struck{$X$}
l'ensemble
\[
  \operatorname{D\acute{e}p}(X) = \{ X'^{\circ} \mid X' \in \widetilde{X} \}
  \subset \Sigma\mathrm{cons} \subset \Sigma
\]
(« \emph{déploiement} » \struck{de $X$} (« spatial ») de $X$). On pourrait le
définir pour tout $F \in \mathfrak{F}(\mathcal{A})$, mais pour l'instant on ne
s'y intéresse que pour les $X \in \mathcal{A}$. On a :

$\operatorname{D\acute{e}p}(X)$ admet un plus grand élément
\add{$\neq \varnothing_{\Sigma}$}, $X^{\circ}$, et est formé d'éléments
\uncertain{mutuellement} disjoints ($\mathrel{|\circ|}$) de
$\Sigma_{\ill{}}$. De plus, l'application
$\widetilde{X} \to \operatorname{D\acute{e}p}(X)$ est bijective, et
$\operatorname{D\acute{e}p}(X)$ hérite de $\widetilde{X}$ une structure
d'ordre. On considérera \uncertain{aussi} $\operatorname{D\acute{e}p}(X)$,
comme une partie de $\Sigma$, munie \struck{de la} \add{de la} relation
d'ordre \ill{}, comme une \emph{figure} (élémentaire) \uncertain{ensembliste}
dans $\Sigma$, « \emph{\uncertain{spatiale}} » i.e. \ill{} \ill{}
\uncertain{spatiale}.

\page{44}
\note{p. 41 de l'auteur.}

Supposons maintenant qu'on ait
\[
  X \ll Y \qquad (X, Y \in \mathcal{A})
\]
--- on a alors une application croissante canonique
\[
  \varphi : \widetilde{X} \longrightarrow \widetilde{Y}
\]
telle que, pour tout $X' \in \widetilde{X}$, on ait
\[
  X' \overset{\circ}{\ll} \text{\struck{$\varphi(X')$}}\
  Y' \overset{\text{déf}}{=} \varphi(X') .
\]
En termes des déploiements, on peut interpréter $\varphi$ par une application
croissante
\[
  \psi : \operatorname{D\acute{e}p}(X) \longrightarrow \operatorname{D\acute{e}p}(Y)
\]
satisfaisant
\[
  \forall S \in \operatorname{D\acute{e}p} X, \quad S \leq \psi(S)
\]
(où $\leq$ est la relation $\subset$ dans $\Sigma$). Cette application $\psi$ est
uniquement déterminée par cette propriété, car \struck{\uncertain{si on a}}
\add{\uncertain{si}}
\[
  S \leq T, T', \quad \text{\uncertain{avec}} \quad
  T, T' \in \operatorname{D\acute{e}p}(Y)
\]
on a $T = T'$ (et ceci, pour peu seulement que $S$ soit un élément de
$\Sigma$, $\neq \varnothing_{\Sigma}$). En effet, si $T \neq T'$, on a
$T \mathrel{|\circ|} T'$,

\page{45}
\note{p. 42 de l'auteur.}

donc $T \cap T' = \varnothing$, donc $S \leq \varnothing$ donc
$S = \varnothing$.

Cela impose donc une condition \add{\uncertain{nécessaire}}, en termes des
\struck{supports} \add{déploiements}, pour qu'on ait
\[
  X \ll \text{\struck{\ill{}}}\ Y ,
\]
c'est la condition suivante :

\uncertain{R$\forall$} \quad $\forall S \in \operatorname{D\acute{e}p} X$,
$\exists T \in \operatorname{D\acute{e}p} Y$, \uncertain{avec}
\add{(nécessairement unique, soit $T = \psi(S)$)} $S \leq T$. De plus,
l'application
\[
  \psi : \operatorname{D\acute{e}p}(X) \longrightarrow \operatorname{D\acute{e}p}(Y)
\]
est croissante.

De plus, on a $X \overset{\circ}{\ll} Y$ ssi l'application $\psi$ précédente
transforme \add{\uncertain{le}} plus grand élément $X^{\circ}$ en le plus
grand élément $Y^{\circ}$, i.e. si et seulement si on a \struck{$\ll$}
\[
  X^{\circ} \leq Y^{\circ}
\]
(\emph{NB} \ill{} \uncertain{supposé} déjà que $X \ll Y$…).

\emph{Définition.} On dit qu'un \struck{\ill{}} \add{\uncertain{atelier}} est
\emph{spatial}, si la condition nécessaire précédente \uncertain{R$\forall$},
sur $\operatorname{D\acute{e}p} X$ et $\operatorname{D\acute{e}p} Y$, pour
qu'on ait $X \ll Y$

\marginal{Si l'on a $X \trianglelefteq Y$, alors
$\psi : \operatorname{D\acute{e}p} X \to \operatorname{D\acute{e}p} Y$ est
une \uncertain{inclusion} d'ensembles ordonnés
($\operatorname{D\acute{e}p}(X) \subset \operatorname{D\acute{e}p}(Y)$, et
\struck{\ill{}} $\operatorname{D\acute{e}p}(X)$ \uncertain{est} une partie
\ill{} de $\operatorname{D\acute{e}p}(Y)$).}

\page{46}
\note{p. 43 de l'auteur.}

est aussi suffisante.

Le terme « spatial » \struck{signifie} est destiné ici à suggérer que la
description de la relation primitive $\ll$ peut se faire en termes de
$\trianglelefteq$, $\mathrel{|\circ|}$, par référence à des notions des notions
purement « spatiales » (la notion de support, qui \struck{fait le lien avec la}
correspond à la notion « d'espace occupé » par une figure.)

Cela nous amène à la question : étant donnée une structure
\[
  \text{\struck{\ill{}}}\ (\mathcal{A}, \trianglelefteq, \mathrel{|\circ|}) ,
\]
à quelles conditions provient-elle d'une structure d'atelier spatial
\[
  (\mathcal{A}, \trianglelefteq, \overset{\circ}{\ll}, \mathrel{|\circ|}) \ ?
\]

La structure de départ \struck{$(\mathcal{A},$}
$(\mathcal{A}, \trianglelefteq, \mathrel{|\circ|})$ * permet de construire
l'ensemble ordonné $\Sigma$, \uncertain{avec} la structure supplémentaire
$\complement$ (permettant de définir $\mathrel{|\circ|}$ sur $\Sigma$), et
l'application
\[
  X \longmapsto \operatorname{D\acute{e}p}(X) \qquad
  \mathcal{A} \longrightarrow \operatorname{Fig\acute{e}l\,ord}(\Sigma)
\]
\note{« ord » est ajouté au-dessus de « Figél ».}

--- $\operatorname{Fig\acute{e}l\,ord}(\Sigma)$ désigne l'ensemble
\struck{des} des figures élémentaires \add{\uncertain{ens.}} \uncertain{ordonnées}
de $\Sigma$ qui

\marginal{* + l'axiome At 4 qu'il faut \uncertain{supposer} \ill{},
\uncertain{comme} \ill{} conditions (i)}

\page{47}
\note{p. 44 de l'auteur.}

sont compatibles avec la relation \add{(d'ensemble \ill{})}
\add{\uncertain{de support}} \struck{\ill{}} \uncertain{condition} ;
\struck{\ill{}} i.e. telles que \struck{$\forall A \in$} $S \in \mathfrak{F}$

a) $S \in \mathfrak{F} \Longrightarrow S \neq \varnothing_{\Sigma}$

b) $S, T \in \mathfrak{F}$, $S \neq T \Longrightarrow
    S \leq \complement T$ i.e. $T \leq \complement S$.


\struck{La première condition nécessaire, qui exprime la}

La condition
\[
  \operatorname{D\acute{e}p}(X) \overset{\circ}{\ll} \operatorname{D\acute{e}p}(Y)
\]
est une relation de préordre sur $\mathcal{A}$, \uncertain{que je noterai}
\[
  X \overset{\circ}{\ll} Y
\]
La première condition, c'est que cette relation soit une \emph{relation
d'ordre}, i.e. \struck{que}

(ii) Soient $X, Y \in \mathcal{A}$, tels que
$\operatorname{D\acute{e}p} X = \operatorname{D\acute{e}p} Y$, alors $X = Y$.

Ceci fait, on trouve une structure
$(\mathcal{A}, \trianglelefteq, \overset{\circ}{\ll}, \mathrel{|\circ|})$,
satisfaisant \uncertain{déjà} At 1 a) b) c). Je voudrais que la relation $\ll$
qu'elle définit soit celle \struck{de} déduite par $\ll$ de
$\operatorname{Fig\acute{e}l\,ord}(\Sigma)$. \struck{Cela signifie que si
$\operatorname{D\acute{e}p}(X)$} Supposons \struck{donc} \struck{Dé}
\[
  \operatorname{D\acute{e}p}(X) \ll \operatorname{D\acute{e}p}(Y)
\]
donc

\page{48}
\note{p. 45 de l'auteur.}

\[
  \operatorname{D\acute{e}p} X \overset{\circ}{\ll} \mathfrak{F}
  \leq \text{\struck{$\ll$}}\ \operatorname{D\acute{e}p}(Y)
\]
\note{un trait relie $\mathfrak{F}$ à l'annotation placée dessous : « sous-figure
élémentaire de $Y$, définie par $\psi(X^{\circ})$ ».}

on sait que $\mathfrak{F}$ sera de la forme \struck{De}
\[
  \mathfrak{F} = \operatorname{D\acute{e}p}(Y'), \qquad Y' \trianglelefteq Y ,
\]
et on aura alors bien alors, par définition, $X \overset{\circ}{\ll} Y'$, ok.

\emph{At 2} \quad a) immédiat \uncertain{par} (ii)

\quad b) Soit $X \overset{\circ}{\ll} Y$, $X' \trianglelefteq X$, prouvons
$\exists Y' \in \widetilde{Y}$ avec $X' \overset{\circ}{\ll} Y'$. Immédiat.
\note{$X'$ est écrit sous $X$, le signe $\trianglelefteq$ tourné
verticalement entre eux.}

\emph{At 3} \quad On note que la relation $X \mathrel{|\circ|} Y$ se
\struck{transpose par} \add{traduit sur}
$\operatorname{D\acute{e}p} X$, $\operatorname{D\acute{e}p}(Y)$ par
\struck{plus}

si $S$ est le plus grand élément de $\operatorname{D\acute{e}p} X$, $T$
------ $\operatorname{D\acute{e}p} Y$, on a $S \mathrel{|\circ|} T$ dans
$\Sigma$.

Et l'axiome At 3 résulte de la relation \uncertain{évidente} dans $\Sigma$
\[
  S \mathrel{|\circ|} T,\ S' \leq S,\ T' \leq T \Longrightarrow
  S' \mathrel{|\circ|} T'
\]

\emph{At 4} \quad \struck{Évident.} Ne fait pas intervenir
$\overset{\circ}{\ll}$, \struck{déjà connu.} il faut bien sûr l'imposer
d'avance.

\page{49}
\note{p. 46 de l'auteur.}

On trouve donc que la structure d'atelier \emph{spatial} peut se décrire par les
données
\[
  (\mathcal{A}, \trianglelefteq, \mathrel{|\circ|})
\]
avec les seuls axiomes

\emph{Atspat 1} \quad a) $\trianglelefteq$ est une relation d'ordre.

\quad b) $\mathrel{|\circ|}$ est symétrique \uncertain{antiréflexive}.

\emph{Atspat 2} \quad Si $X \trianglelefteq Y$, on a
$X \mathrel{|\circ|} Y$.
\note{ainsi sur la page ; une barre de négation sur la relation n'est pas
exclue.}

\emph{Atspat 3} \quad L'application
\[
  \mathcal{A} \longrightarrow
  \text{\struck{$\mathfrak{F}$}}\ \operatorname{Fig\acute{e}l}(\mathfrak{P}(\mathcal{A}))
\]
\[
  X \longmapsto \operatorname{D\acute{e}p} X
  = \bigl( \{ X'^{\circ} \mid X' \trianglelefteq X \} ,
\]
\struck{ordre} + relation d'ordre sur $\operatorname{D\acute{e}p} X$ déduite
de celle de
$\widetilde{X} = \{ X' \in \mathcal{A} \mid X' \trianglelefteq X \}$ )
\struck{où $X'^{\circ} = \{ Y$}

est \emph{injective}. (Dans cette définition, \uncertain{pour}
$Y \in \mathcal{A}$,
\[
  Y^{\circ} = \{ Z \in \mathcal{A} \mid \forall Z' \in \mathcal{A},\
  Z' \mathrel{|\circ|} Y \Rightarrow Z' \mathrel{|\circ|} Z \} \ ) .
\]

Je vais essayer de reformuler la structure, avec
\[
  \Sigma = \Sigma_{\mathcal{A}}
\]
comme ensemble de base.

\emph{Définition.} Appelons \emph{spatios} (ensembliste) un
\struck{ensemble ordonné} \add{triple} $(\Sigma, \leq, \complement)$, où
$\Sigma$ est un ensemble, $\leq$ une relation d'ordre sur $\Sigma$,
$\complement$ une application de $\Sigma$ en lui-même,

\page{50}
\note{p. 47 de l'auteur.}

satisfaisant aux axiomes

\emph{Spat 1}. $\leq$ est une relation d'ordre, admettant un plus petit
élément (noté $0_{\Sigma}$) et un plus grand élément (noté $1_{\Sigma}$), et
\uncertain{on admet} des Inf quelconques (donc aussi des Sup quelconques).

\emph{Spat 2} \quad \struck{$\complement$ est une} $\complement^{2} = \mathrm{id}$
($\complement$ est une involution) et \struck{\ill{}}.
\[
  x \leq y \iff \complement y \leq \complement x
\]
(i.e. $\complement$ est un anti-homomorphisme involutif de l'ens. ordonné sur
lui-même).

\emph{Spat 3} \quad $\forall x$, on a
$x \wedge \complement x = 0_{\Sigma}$ (ce qui équivaut, par $\complement$, à :
$\complement x \vee x = 1_{\Sigma}$)

Dans un tel ensemble, on met
\[
  X \mathrel{|\circ|} Y \iff X \leq \complement Y
  \quad (\iff Y \leq \complement X) .
\]
On a
\[
  X \mathrel{|\circ|} Y \Longrightarrow X \wedge \text{\struck{\ill{}}}\ Y = 0_{\Sigma}
\]
On trouve
\[
  X \mathrel{|\circ|} Y \ \text{est relation symétrique}
\]
\[
  X \mathrel{|\circ|} X \iff X = 0_{\Sigma}
\]
donc dans $\Sigma^{*} = \Sigma \setminus \{ 0_{\Sigma} \}$, la relation
$\mathrel{|\circ|}$ est antiréflexive, en plus d'être symétrique.

\emph{Ex} \quad Soit $\mathcal{A}$ un ens. muni d'une relation
$\mathrel{|\circ|}$ qui est symétrique, soit $\mathcal{A}_{0}$ l'ens. des
$X \in \mathcal{A}$ tels que $X \mathrel{|\circ|} X$ \struck{\ill{}}
\struck{supposons que $X \in \mathcal{A}_{0}$}
($\mathcal{A}_{0} = \varnothing$ dans le cas où $\mathcal{A}$ est un atelier,
et $\mathrel{|\circ|}$ sa relation de disjonction \uncertain{intrinsèque}…),
\struck{\ill{}}
\struck{supposons que} $X \in \mathcal{A}_{0}$ $\iff$
$X \mathrel{|\circ|} Y$ $\forall Y \in \mathcal{A}$. On trouve alors une notion
de support et de \uncertain{cosupport}, et un ensemble \struck{\ill{}} de
\[
  \Sigma \subset \mathfrak{P}(\mathcal{A})
\]

\page{51}
\note{p. 48 de l'auteur.}

de supports, stable dans $\mathfrak{P}(\mathcal{A})$ par $\bigcap$,
\uncertain{quelconques}, et donc \uncertain{aussi} par des Inf quelconques,
donc aussi des Sup quelconques. Le plus petit élément est
$\mathcal{A}_{00}$
\add{$= \{ X \in \mathcal{A} \mid X \mathrel{|\circ|} Y\ \forall Y \in \mathcal{A} \}$},
le plus grand est $\mathcal{A}$. On a une opération $\complement$
\uncertain{ici} ($\complement^{2} = \mathrm{id}$ (symétrie de
$\mathrel{|\circ|}$). Ainsi, les conditions \struck{a) b) c)}
\add{Spat 1, 2} d'un spatios sont satisfaites. On a \struck{\ill{}}
\add{Spat 3} ssi
\[
  X \mathrel{|\circ|} X \Longrightarrow X \in \mathcal{A}_{00}
\]
--- c'est la condition que j'avais mise dès le début, \uncertain{on dirait}.

Appliquons cette construction au cas
$\mathcal{A} = \Sigma$, $\mathrel{|\circ|}_{\Sigma}$ ! On trouve
\uncertain{\ill{}}
\[
  \Sigma_{0} = \{ 0_{\Sigma} \} = \Sigma_{00} .
\]
Soit \struck{\ill{}} $A \subset \Sigma$, son cosupport $A$ est \uncertain{par
définition}
\[
  \operatorname{cosupp}(A) = \{ Z \in \Sigma \mid \forall X \in A,\
  X \mathrel{|\circ|} Z \}
\]
\note{sous $X \mathrel{|\circ|} Z$ : « i.e. $Z \leq \complement X$ ».}
\[
  = \Bigl\{ Z \in \Sigma ,\ Z \leq \inf_{X \in A} \complement X
  = \complement \operatorname{Sup} A \Bigr\}
\]
\[
  = \widetilde{\complement \operatorname{Sup} A} .
\]
\note{en marge à droite : « $= \complement \operatorname{Sup} A$ ».}

Donc $\operatorname{cosupp} A$ ne dépend que de $\operatorname{Sup} A$
\add{($= \xi$, \uncertain{disons})}. On a donc
\[
  \operatorname{supp} A \overset{\text{déf}}{=}
  \operatorname{cosupp} \operatorname{cosupp} A
  = \complement\bigl(\operatorname{Sup} \widetilde{\complement(\xi)}\bigr)
  = \widetilde{\xi}
\]
\note{sous $\operatorname{Sup} \widetilde{\complement(\xi)}$, une accolade :
« $\complement(\xi)$ ».}

De cette façon, on voit que l'ensemble des supports s'identifie à
\uncertain{$\Sigma$} via $\Sigma \to \widetilde{\Sigma}$,
\note{le dernier signe, un $\Sigma$ surmonté d'un trait ou d'un tilde, est
lu sous réserve.}

\marginal{Si on \uncertain{prend} $\mathcal{A} = \Sigma^{*} =
\Sigma \setminus \{ 0_{\Sigma} \}$, on trouve \uncertain{encore} \ill{} de
$\Sigma$, \ill{} \uncertain{cette fois} d'une relation $\mathcal{A}$
\uncertain{symétrique} \ill{}. On a $\mathcal{A}_{00} = \varnothing$ \ill{}.
* \ill{} : \ill{} \uncertain{vide}}

\page{52}
\note{p. 49 de l'auteur.}

et via cette identification, le support d'une partie \add{$A$} de $\Sigma$
s'identifie à \uncertain{son} Sup $\xi$, son cosupport à $\complement \xi$.

\struck{Ceci dit, à un atelier spatial $(\mathcal{A}, \trianglelefteq,
\mathrel{|\circ|})$ est associé un spatios $(\Sigma, \leq, \complement)$,
\ill{} en plus \ill{}}
\note{ce paragraphe est barré de traits obliques.}

Étant donné un spatios $\Sigma$, on \struck{définit} appelle \emph{figure du
spatios} \struck{$\operatorname{Fig\acute{e}l}(\Sigma)$}
\struck{$\operatorname{Figspat}(\Sigma)$}
\struck{$= \{ (\mathfrak{A}, \rho) \mid$} \struck{$\mathfrak{F} \in$} un \struck{partie} couple $(\mathfrak{F}, \rho)$, où
\[
  \mathfrak{F} \subset \Sigma
\]
est une partie de $\Sigma$, et $\rho$ une relation sur $\mathfrak{F}$,
satisfaisant les conditions

\emph{Figspat 1} \quad $\rho$ est une relation d'ordre sur $\mathfrak{F}$

\emph{Figspat 2} \quad a) $\forall X \in \mathfrak{F}$, on a
$X \neq 0_{\Sigma}$, et

\quad b) $\forall X, Y \in \mathfrak{F}$, $X \neq Y \Longrightarrow
X \mathrel{|\circ|} Y$ (i.e. $X \leq \complement Y$)

\emph{NB} On pourrait être plus exigeant, en s'inspirant de l'axiome des
supports pour un atelier. On ne le fait pas pour le moment.

On désigne par $\operatorname{Figspat}(\Sigma)$ l'ensemble des figures
spatiales dans $\Sigma$, et par

\page{53}
\note{p. 50 de l'auteur.}

\struck{$\operatorname{Figspat\acute{e}l}(\Sigma)$}
\add{$\operatorname{Atspat}(\Sigma)$}, l'ensemble des figures spatiales
\emph{élémentaires} (i.e. telles que $\mathfrak{F}$ ait un plus grand
élément). On va mettre une structure d'atelier sur
$\operatorname{At}\ill{}(\Sigma)$, ainsi :
\[
  \text{\struck{$(\mathfrak{F}, \rho) \trianglelefteq \mathfrak{F}'$}} \qquad
  \operatorname{Atspat}(\Sigma) \subset \operatorname{Atens}(\Sigma)
\]

1) La relation $\trianglelefteq$ dans $\operatorname{Atspat}(\Sigma)$ est
induite par la relation similaire dans $\operatorname{Atens}(\Sigma)$.

2) La relation $\overset{\circ}{\ll}$ dans $\operatorname{Atspat} \Sigma$ est
\add{définie} \struck{induite} plus haut :
$(\mathfrak{F}, \rho) \overset{\circ}{\ll} (\mathfrak{F}', \rho')$
\add{ssi} \struck{par la relation similaire dans $\operatorname{Atens}
\Sigma$} $\forall X \in \mathfrak{F}$, $\exists X' = \psi(X) \in
\mathfrak{F}'$ \uncertain{avec} $X \leq$ \struck{$X'$} $\psi(X)$ (NB.
\struck{\ill{}} cet $X'$ est unique), et
$\psi : \mathfrak{F} \to \mathfrak{F}'$ est croissante \struck{\ill{}} et
transforme plus grand élément en \uncertain{plus grand}.

3) La relation $\mathrel{|\circ|}$ dans $\operatorname{Atspat}(\Sigma)$ est
\struck{\ill{}} plus stricte que dans $\operatorname{Atens}(\Sigma)$ : on pose
$(\mathfrak{F}, \rho) \mathrel{|\circ|} (\mathfrak{F}', \rho')$, ssi (posant
$X$ = plus grand élément de $\mathfrak{F}$, de même pour $X'$) non seulement on
a $X \neq X'$ (relation $\mathrel{|\circ|}$ dans $\operatorname{Atens}$) mais
même $X \mathrel{|\circ|} X'$ (ce qui est plus fort, vu que $X, X' \neq
0_{\Sigma}$).

Je dis que pour ces relations, on a \emph{un atelier}, \struck{spatial} et
même \ill{} un \emph{atelier spatial}.

\page{54}
\note{p. 51 de l'auteur.}

\emph{\struck{Atsp} At 1} \quad Clair

\emph{At 2} \quad Clair \struck{(c'est l'axiome ensembliste)}

\emph{At 3} \quad Clair

\emph{At 4} \quad Clair.

\struck{\ill{}} Voyons quel est le \struck{\ill{} des} \struck{\ill{}}
spatios associé à précédent atelier. \struck{\ill{}} Soit
$\mathcal{A} = \operatorname{Atspat}(\Sigma)$. Considérons
\struck{$(\mathfrak{F}, \rho)$} l'application
\[
  \mathcal{A} \longrightarrow \Sigma
\]
\[
  X = (\mathfrak{F}, \rho) \longmapsto \text{plus grand élément de } \mathfrak{F},
\]
\add{(\uncertain{noté} $X^{\circ}$, cf justification \uncertain{plus bas})}

Ceci dit, la relation $\mathrel{|\circ|}$ de $\mathcal{A}$ est induite, via
cette application, par la relation \struck{\ill{}} de \uncertain{même nom}
dans $\Sigma$. On trouve donc une application injective
\[
  \Sigma_{\mathcal{A}} \longrightarrow \Sigma_{\Sigma} = \Sigma .
\]
\note{sous $\Sigma_{\Sigma}$, deux mots : \ill{} \ill{}.}

Mais \uncertain{on a} aussi
\[
  \Sigma \longrightarrow \mathcal{A}
\]
\[
  X \longmapsto \text{\struck{$\{X\}$}}\ (\{X\}, \rho_{X}) \qquad
  \rho_{X} \ \text{la relation d'ordre (unique) sur } \{X\}
\]
\note{« (unique) » est ajouté au-dessus de la ligne.}
telles que la relation $\mathrel{|\circ|}$ sur $\Sigma$ soit induite par celle
sur $\mathcal{A}$, et le composé
\[
  \Sigma \longrightarrow \mathcal{A} \longrightarrow \Sigma
\]
est l'identité.

\page{55}
\note{p. 52 de l'auteur.}

d'où
\[
  \Sigma_{\mathcal{A}} \simeq \Sigma_{\Sigma} \simeq \Sigma
\]
et on voit que \uncertain{ça} \uncertain{respecte} relations d'ordre et
$\complement$.

Ceci dit, l'application canonique
\[
  \mathcal{A} \longrightarrow \operatorname{Atspat}(\Sigma_{\mathcal{A}})
\]
s'identifie, via $\Sigma_{\mathcal{A}} \simeq \Sigma$, en composant
\[
  \mathcal{A} \longrightarrow \operatorname{Atspat}(\Sigma_{\mathcal{A}})
  \xrightarrow{\ \sim\ } \operatorname{Atspat}(\Sigma) = \mathcal{A}
\]
à l'application identique. Cette application \ill{} \ill{} \uncertain{non
seulement} injection, mais \emph{bijection} !

Revenons au cas d'un atelier \struck{spatial} quelconque $\mathcal{A}$.
\add{Posons $\Sigma = \Sigma_{\mathcal{A}}$.} Alors \add{\uncertain{on a un}
\emph{homom. d'ateliers}} \struck{$\mathcal{A}$ s'identifie comme un \ill{}
atelier de $\operatorname{Atspat}(\Sigma)$}
\[
  \mathcal{A} \longrightarrow \operatorname{Atspat}(\Sigma)
\]
\[
  X \longmapsto \operatorname{D\acute{e}p}(X) ,
\]
(même si \struck{$X$} $\mathcal{A}$ \uncertain{n'est} pas spatial). C'est un
\uncertain{homom.} \emph{fidèle}, i.e.
\[
  X \mathrel{|\circ|} Y \iff
  \operatorname{D\acute{e}p}(X) \mathrel{|\circ|} \operatorname{D\acute{e}p} Y .
\]
Dire que $\mathcal{A}$ est spatial, signifie

\page{56}
\note{p. 53 de l'auteur.}

qu'on a un homom. injectif, qui identifie \add{\uncertain{donc}}
$\mathcal{A}$ à un sous-atelier de $\operatorname{Atspat}(\Sigma)$,
\struck{soit $\mathcal{A}$}.

Ainsi, on trouve, associé à $\mathcal{A}$,

1°) un spatios $\Sigma$

2°) un sous-atelier $\mathcal{A} \subset \operatorname{Atspat}(\Sigma)$.
\struck{i.e.}

Sous-atelier signifie ici, simplement, une \emph{partie} \add{de Atspat}, qui
satisfait la condition Sousat (cf p. 27) --- mais en fait, c'est \struck{à}
une partie \emph{fermée} pour $\trianglelefteq$. Cette condition joue le rôle
de la \add{\uncertain{p. 9}} condition Atens 1 (cf \uncertain{Ch.}
\emph{VIII}) pour les ateliers ensemblistes, rapportés à
\add{(ens. de base $L$)} $L$. Il faut une condition qui \struck{\ill{}}
tienne lieu de Atens 2 ($\forall x \in L$, $\{x\} \in \mathcal{A}$). Ce
serait

Atspat 2. Ci-dessous.

\note{« p. 27 » et « p. 9 » renvoient à la pagination de l'auteur, hors de ce
lot.}

\struck{Donc} On va donc reprendre la structure sur l'ens. de base $\Sigma$,
\uncertain{avec} comme données
($\leq$, \struck{\ill{}} $\complement$, $\mathcal{A}$) \struck{\ill{}},
où $\leq$, $\complement$ font de $\Sigma$ un spatios, et où $\mathcal{A}$ est
une \struck{partie} ens. de figures spatiales élémentaires
\uncertain{relativement} à $\Sigma$
\[
  \mathcal{A} \subset \operatorname{Atspat}(\Sigma) ,
\]

\page{57}
\note{p. 54 de l'auteur.}

satisfaisant les deux seules (?) conditions

\emph{Atspat 1'} \quad Si $X \in \mathcal{A}$, et
$Y \trianglelefteq X$, alors $Y \in \mathcal{A}$ ($\mathcal{A}$ fermé dans
$\operatorname{Atspat}(\Sigma)$, pour $\trianglelefteq$).

\struck{\emph{Atspat 2} \quad $\forall S \in \Sigma$,} \struck{\ill{}
soit} \struck{$A_{S} = \{ X \in \mathcal{A} \mid$
$X^{\circ} \leq S \}$} \struck{\ill{} $\forall Y \in \widetilde{X}$, \ill{}}
\struck{$\Sigma_{0} = \{ S \in \Sigma \mid$ $\exists X \in \mathcal{A}$}
\struck{$+\ X^{\circ} \mathrel{|\circ|} X$}

\emph{Atspat 2'} \quad \struck{$\forall$} Soit
\[
  \Sigma' = \{ X^{\circ} \mid X \in \mathcal{A} \} \subset \Sigma .
\]
Pour tout $S \in \Sigma$ soit
$\Sigma'_{\leq S} = \{ \xi \in \Sigma' \mid \xi \leq S \}$. Ceci posé, on a,
pour tout $S \in \Sigma$
\[
  S = \operatorname{Sup} \Sigma'_{\leq S} .
\]

Je m'attends que, pour de telles données, \struck{l'atelier $\mathcal{A}$}
\uncertain{partie} de $\operatorname{Atspat}(\Sigma)$, muni de la structure
d'atelier induite, est un atelier spatial, et que l'on a un iso. can.
\[
  \Sigma \simeq \Sigma_{\mathcal{A}}
\]
moyennant \uncertain{lequel} l'application canonique
\[
  \mathcal{A} \xrightarrow{\ \operatorname{D\acute{e}p}\ }
  \operatorname{Atspat}(\Sigma_{\mathcal{A}})
\]
devient l'\emph{inclusion} de $\mathcal{A}$ dans
$\operatorname{Atspat}(\Sigma)$.

\page{58}
\note{p. 55 de l'auteur.}

\marginal{\ill{} \uncertain{justifié}}

Il sera plus commode de vérifier, \add{d'abord,} en termes des données
$\trianglelefteq_{\mathcal{A}}$, $\mathrel{|\circ|}_{\mathcal{A}}$ induites par
$\operatorname{Fig\acute{e}lspat}(\Sigma)$, les axiomes Atspat 1--3.
Atspat 1, 2 triviales. Pour Atspat 3 on va déterminer \add{\uncertain{le
spatios}} $\Sigma_{\mathcal{A}}$ (défini en termes de
$\mathrel{|\circ|}_{\mathcal{A}}$ seulement), et l'application
\[
  \mathcal{A} \xrightarrow{\ \operatorname{D\acute{e}p}\ }
  \operatorname{Figspat\acute{e}l}(\Sigma_{\mathcal{A}})
\]
(défini en termes de $\mathrel{|\circ|}_{\mathcal{A}}$ et
$\trianglelefteq_{\mathcal{A}}$), et voir que \uncertain{ces} s'identifient
resp. à $\Sigma$ et à l'inclusion
$\mathcal{A} \subset \operatorname{Figspat\acute{e}l}(\Sigma_{\mathcal{A}})$.

\struck{Soit} Considérons les applications
\[
  \mathcal{A} \xrightarrow{\ \alpha\ } \Sigma
  \xrightarrow{\ \beta\ } \mathfrak{P}(\mathcal{A})
\]
\[
  S \longmapsto \{ X \in \mathcal{A} \mid X^{\circ} \leq S \}
\]
\[
  X = (\mathfrak{F}, \rho) \longmapsto X^{\circ} =
  \text{plus grand él. de } \mathfrak{F} \text{ pour } \rho
\]
On a
\[
  \beta(\alpha(X)) = \{ Z \in \mathcal{A} \mid Z^{\circ} \leq X^{\circ} \}
  \subset \operatorname{supp}^{\circ}_{\mathcal{A}} X
\]
\struck{$\operatorname{supp}^{\circ} X$} et
\[
  \alpha\beta(S) = \{ X^{\circ} \mid X \in \mathcal{A},\ \text{t.q. }
  X^{\circ} \leq S \} \subset \operatorname{supp}^{\circ}_{\Sigma}(S)
\]
D'où résulte que $\alpha$ et $\beta$ induisent des bijections inverses l'une
de l'autre entre
\[
  \Sigma_{\mathcal{A}} \quad \text{et} \quad \Sigma_{\Sigma} .
\]
Utilisant
\[
  \Sigma_{\alpha} : \Sigma_{\mathcal{A}} \longrightarrow \Sigma_{\Sigma}
  \overset{\operatorname{supp}^{\circ}_{\Sigma}}{\rightleftarrows} \Sigma ,
\]
\note{sous la dernière flèche : « cf plus haut p. 48 » --- p. 48 de l'auteur,
c'est-à-dire la page 51 de ce lot.}

\struck{on trouve l'application} pour identifier $\Sigma_{\mathcal{A}}$ et
$\Sigma$, et interprétons

\marginal{Cf suite plus explicite \ill{} \ill{} \uncertain{l'associativité}
\ill{} de $\Sigma$…}

\page{59}
\note{p. 56 de l'auteur.}

en conséquence $\operatorname{D\acute{e}p}_{\mathcal{A}}$ comme
\[
  \mathcal{A} \xrightarrow{\ \operatorname{D\acute{e}p}^{!}_{\mathcal{A}}\ }
  \text{\struck{\ill{}}}\ \mathfrak{P}(\Sigma_{\mathcal{A}})
  \longrightarrow \mathfrak{P}(\Sigma)
\]
\[
  X \longmapsto \{ \operatorname{supp}^{\circ}_{\mathcal{A}}(X') \mid
  X' \in \widetilde{X} \} \overset{?}{=}
\]
\uncertain{facile}
\note{sous $X$ : « $= (\mathfrak{F}, \rho)$ » ; sous $\widetilde{X}$ :
« $= \mathfrak{F}$ ». Un premier but, raturé, est surchargé de
$\mathfrak{P}(\Sigma_{\mathcal{A}})$.}

\marginal{(Cas pour \uncertain{l'instant} \uncertain{prenant}
$\operatorname{D\acute{e}p}_{\mathcal{A}} X$ \uncertain{sans} la relation
d'ordre, soit $\operatorname{D\acute{e}p}^{!}_{\mathcal{A}} X$)}

Or $\operatorname{supp}^{\circ}_{\mathcal{A}}(X')$
\add{($\in \Sigma_{\mathcal{A}}$)} devient, pour $X' \in \mathfrak{F}$
\add{(donné $\subset \Sigma$)}, par
$\Sigma_{\mathcal{A}} \xrightarrow{\ \sim\ } \Sigma_{\Sigma}$,
$\operatorname{supp}^{\circ}_{\Sigma}(X')$, et par
$\Sigma_{\Sigma} \longleftarrow \Sigma$, $X'$. Donc l'image de
$X = (\mathfrak{F}, \rho)$ \struck{par} \uncertain{par le} composé (*)
\uncertain{n'est autre} que $\mathfrak{F}$, et la relation d'ordre
\uncertain{induite} dessus est celle déduite par transport de structure de
celle de $\widetilde{X} = \mathfrak{F}$, par l'application identique de
$\widetilde{X} \xrightarrow{\ \sim\ } \mathfrak{F}$, i.e. c'est $\rho$.

On \uncertain{gagne} !

\emph{Scholie} \quad La donnée d'un atelier spatial revient à celle d'un
spatios $(\Sigma, \leq, \complement)$ \add{(\uncertain{ou}
$\overset{\circ}{\ll}$)}, muni de la structure supplémentaire définie par une
$\mathcal{A} \subset \operatorname{Figspat\acute{e}l}(\Sigma)$, satisfaisant
les deux conditions Atspat 1', 2' p. 54. Ainsi l'ensemble de \uncertain{base} de
la

\note{« p. 54 » de l'auteur : page 57 de ce lot.}

\page{60}
\note{p. 57 de l'auteur.}

structure \uncertain{envisagée} (qui avait été, tour à tour, $\mathfrak{F}$,
(mais \ill{} \uncertain{moment} !), $\mathcal{A}$ ou
$(\mathcal{A}, \mathfrak{F})$, $\mathcal{L}$, est à présent $\Sigma$, le
« spatios des supports ». [Ce dernier est particulièrement commode pour
formuler les axiomes des supports, cf plus bas, même pour des ateliers
\uncertain{non nécessairement} spatiaux.]

\emph{Sorite} du \struck{\ill{}} spatios \struck{$\Sigma(E)$}
\uncertain{associé à} un ens. \struck{$E$} muni d'une
\add{\uncertain{disjonction}} relation symétrique \add{anti}réflexive
(\struck{\ill{}} \uncertain{notée} $\mathrel{|\circ|}_{X}$), --- \uncertain{ou
plus généralement} \ill{}, d'une relation telle que
\[
  X \mathrel{|\circ|} X \Longrightarrow \forall Y \in E,\
  X \mathrel{|\circ|} Y .
\]

1°) Soient $E = (E, \mathrel{|\circ|}_{E})$,
$E' = (E', \mathrel{|\circ|}_{E'})$ munis de relations symétriques, une
correspondance
\[
  E \xrightarrow{\ \varphi\ } \mathfrak{P}(E')
\]
(\uncertain{non nécessairement} compatible à $\mathrel{|\circ|}_{E}$,
$\mathrel{|\circ|}_{E'}$) est dite un \emph{morphisme} de $E$ dans $E'$, si
\struck{$\forall x, y \in E$,}
\[
  \forall x, y \in E, \ \text{et } x' \in \varphi(x),\ y' \in \varphi(y),
  \quad x \mathrel{|\circ|} y \Longrightarrow x' \mathrel{|\circ|} y'
\]
Pour cette notion de morphisme, et la \struck{notion de} composition des
correspondances,

\marginal{\emph{NB} La \uncertain{catégorie} \ill{} \ill{} des \ill{} :
celles des \ill{} $\forall y \neq x$ \ill{} ; $x \mathrel{|\circ|} x$
\ill{} \ill{} \uncertain{ces points}}

\marginal{\emph{Ex.} La corr. \uncertain{identique} de $E$ \ill{}}

\note{La phrase et le sorite se poursuivent au-delà de ce lot.}

\end{document}
